\begin{definition}[The Lie hull of a frame of vector fields]
Let $n,m\in\mathbb N$ and let
$\sigma:\mathbb R^n\to\mathbb R^{n\times m}$, $\sigma(x)=[\sigma^1(x),\ldots,\sigma^m(x)]$,
be a frame of $m$ vector fields on $\mathbb R^n$. Call a family $S$ of vector fields on
$\mathbb R^n$ \emph{admissible for $\sigma$} if it contains each $\sigma^i$ and is closed
under the Lie bracket, i.e. $[V,W]\in S$ whenever $V,W\in S$. The \emph{Lie hull}
$\mathrm{LieHull}\ \sigma$ is the intersection of all admissible families:
\[
 \mathrm{LieHull}\ \sigma\ =\ \bigcap\big\{\,S\ :\ S\text{ admissible for }\sigma\,\big\}.
\]

The intersection of admissible families is admissible, so $\mathrm{LieHull}\ \sigma$ is the
smallest family of vector fields containing $\sigma^1,\ldots,\sigma^m$ and closed under
brackets: it consists precisely of the iterated brackets of the frame fields. Its purpose
is to express H\"ormander's rank condition (1.6),
\[
 \dim\Big(\mathrm{Lie}\big(X_1(x),\ldots,X_m(x)\big)\Big)=n\qquad\text{for all }x\in\mathbb R^n,
\]
as the statement that the span of $\{V(x):V\in\mathrm{LieHull}\ \sigma\}$ is all of
$\mathbb R^n$ for every $x$. Taking the smallest such family, rather than merely asserting
the existence of \emph{some} bracket-closed family whose values span, is what makes the
condition a restriction: the family of all vector fields is bracket-closed and its values
span trivially.
\end{definition}
